\documentclass[10pt,a4paper]{amsart}
\usepackage{amsmath,amssymb,amsthm,graphicx}
\usepackage{verbatim}
\newenvironment{reference}{\comment}{\endcomment}
%\newenvironment{reference}{}{}
\newenvironment{slogan}{\comment}{\endcomment}
\newenvironment{history}{\comment}{\endcomment}

% For commutative diagrams we use Xy-pic
\usepackage[all]{xy}

% We use 2cell for 2-commutative diagrams.
\xyoption{2cell}
% Theorem environments.
%
\theoremstyle{plain}
\newtheorem{theorem}[subsection]{Theorem}
\newtheorem{proposition}[subsection]{Proposition}
\newtheorem{lemma}[subsection]{Lemma}

\theoremstyle{definition}
\newtheorem{definition}[subsection]{Definition}
\newtheorem{example}[subsection]{Example}
\newtheorem{exercise}[subsection]{Exercise}
\newtheorem{situation}[subsection]{Situation}

\theoremstyle{remark}
\newtheorem{remark}[subsection]{Remark}
\newtheorem{remarks}[subsection]{Remarks}

\numberwithin{equation}{subsection}

% Macros
%
\def\lim{\mathop{\mathrm{lim}}\nolimits}
\def\colim{\mathop{\mathrm{colim}}\nolimits}
\def\Spec{\mathop{\mathrm{Spec}}}
\def\Hom{\mathop{\mathrm{Hom}}\nolimits}
\def\Ext{\mathop{\mathrm{Ext}}\nolimits}
\def\SheafHom{\mathop{\mathcal{H}\!\mathit{om}}\nolimits}
\def\SheafExt{\mathop{\mathcal{E}\!\mathit{xt}}\nolimits}
\def\Sch{\mathit{Sch}}
\def\Mor{\mathop{\mathrm{Mor}}\nolimits}
\def\Ob{\mathop{\mathrm{Ob}}\nolimits}
\def\Sh{\mathop{\mathit{Sh}}\nolimits}
\def\NL{\mathop{N\!L}\nolimits}
\def\CH{\mathop{\mathrm{CH}}\nolimits}
\def\proetale{{pro\text{-}\acute{e}tale}}
\def\etale{{\acute{e}tale}}
\def\QCoh{\mathit{QCoh}}
\def\Ker{\mathop{\mathrm{Ker}}}
\def\Im{\mathop{\mathrm{Im}}}
\def\Coker{\mathop{\mathrm{Coker}}}
\def\Coim{\mathop{\mathrm{Coim}}}

% Boxtimes
%
\DeclareMathSymbol{\boxtimes}{\mathbin}{AMSa}{"02}

%
% Macros for moduli stacks/spaces
%
\def\QCohstack{\mathcal{QC}\!\mathit{oh}}
\def\Cohstack{\mathcal{C}\!\mathit{oh}}
\def\Spacesstack{\mathcal{S}\!\mathit{paces}}
\def\Quotfunctor{\mathrm{Quot}}
\def\Hilbfunctor{\mathrm{Hilb}}
\def\Curvesstack{\mathcal{C}\!\mathit{urves}}
\def\Polarizedstack{\mathcal{P}\!\mathit{olarized}}
\def\Complexesstack{\mathcal{C}\!\mathit{omplexes}}
% \Pic is the operator that assigns to X its picard group, usage \Pic(X)
% \Picardstack_{X/B} denotes the Picard stack of X over B
% \Picardfunctor_{X/B} denotes the Picard functor of X over B
\def\Pic{\mathop{\mathrm{Pic}}\nolimits}
\def\Picardstack{\mathcal{P}\!\mathit{ic}}
\def\Picardfunctor{\mathrm{Pic}}
\def\Deformationcategory{\mathcal{D}\!\mathit{ef}}
\usepackage{cleveref}
\crefname{originalreference}{original reference}{original references}
\Crefname{originalreference}{Original reference}{Original references}
\makeatletter
\newlabel{lemma-constructible-serre}{{09B5}{source}{}{sourcepointer-1}{}}
\newlabel{lemma-constructible-serre@cref}{{[originalreference][0][]09B5}{[1][0][]source}{}{}{}}
\makeatother
\crefname{originalreference}{original reference}{original references}
\Crefname{originalreference}{Original reference}{Original references}
\makeatletter
\newlabel{lemma-tensor-c}{{0961}{source}{}{sourcepointer-2}{}}
\newlabel{lemma-tensor-c@cref}{{[originalreference][0][]0961}{[1][0][]source}{}{}{}}
\makeatother
\title{Constructibility modulo an ideal propagates from the first reduction to all powers}
\author{The Stacks Project Authors}
\date{stacks-project:proetale.tex:local-snapshot-20261012; GNU FDL 1.2 or later}
\begin{document}
\maketitle
\noindent\textit{Editorial note.} All delivered mathematical text is the author's own, copied from the pinned Stacks Project chapter file proetale.tex; no statement or proof step was rewritten, reordered or removed.\par\medskip
\noindent\textit{Editorial note.} The counted proof is the author's complete argument: the distinguished triangle comparing successive reductions, the identification of its third term as a tensor product of the first reduction together with the resulting constructibility, and the induction using the Serre subcategory properties.\par\medskip
\noindent\textit{Editorial note.} The label lemma-compare-truncations-constructible names the counted statement, and the proof environment immediately following it in the source is delivered verbatim as its proof.\par\medskip
\noindent\textit{Editorial note.} The counted statement is carried by the author's own lemma environment in the chapter source.\par\medskip
\noindent\textit{Editorial note.} Pre-open gate: python operations/agent-authoring/tools/precard\_gate.py --source proetale.tex --label lemma-compare-truncations-constructible --start 5620 --end 5648 --local-id 28003 returned RESULT=OK (L28003 free, occupancy FREE, public interval clear). Candidate taken from sources/stacks-proetale/candidates-halfA.json.\par\medskip
\noindent\textit{Editorial note.} The statement and proof are copied from the pinned chapter file; the cited inputs are kept as external pointers with their canonical Stacks tags.\par\medskip
\noindent\textit{Allowed standard prerequisites (editorial).}\begin{itemize}
\item The distinguished triangle comparing successive reductions. Standard consequences of the octahedral axiom, giving the distinguished triangle K tensor I\textasciicircum{}n/I\textasciicircum{}\{n+1\} -> K\_\{n+1\} -> K\_n -> K tensor I\textasciicircum{}n/I\textasciicircum{}\{n+1\}[1]. Supplies the triangle along which the induction on n proceeds.
\item Constructibility of the tensor product with a quotient of powers of the ideal. Same chapter, the lemma showing that tensoring a constructible complex over Lambda/I with the constant sheaf on I\textasciicircum{}n/I\textasciicircum{}\{n+1\} stays constructible. Supplies the constructibility of the cohomology of the tensor product appearing in the triangle.
\item Stability of constructible complexes along distinguished triangles. Same chapter, the lemma that the class of constructible complexes is a Serre subcategory stable under the operations of triangulated categories. Supplies the induction step, since constructibility of the outer terms of the triangle gives it for the middle one.
\end{itemize}\medskip
\ifdefined\hypertarget\hypertarget{sourcepointer-1}{}\fi
\noindent\textit{Source reference.} Original reference 09B5: sources/stacks-proetale/pinned-original-source.tex, source lines 5560--5560. Prerequisite pointer invoked at proof line(s) 5646; it names a standard fact of the same chapter that the counted argument uses, and every step of the counted proof built on it is retained verbatim in the delivered body.\par\medskip
\ifdefined\hypertarget\hypertarget{sourcepointer-2}{}\fi
\noindent\textit{Source reference.} Original reference 0961: sources/stacks-proetale/pinned-original-source.tex, source lines 5605--5605. Prerequisite pointer invoked at proof line(s) 5643; it names a standard fact of the same chapter that the counted argument uses, and every step of the counted proof built on it is retained verbatim in the delivered body.\par\medskip
\begin{lemma}
\label{lemma-compare-truncations-constructible}
Let $X$ be a scheme. Let $\Lambda$ be a Noetherian ring.
Let $I \subset \Lambda$ be an ideal.
Let $K$ be an object of $D(X_\proetale, \Lambda)$.
Set $K_n = K \otimes_\Lambda^\mathbf{L} \underline{\Lambda/I^n}$.
If $K_1$ is in $D^-_c(X_\proetale, \Lambda/I)$, then
$K_n$ is in $D^-_c(X_\proetale, \Lambda/I^n)$ for all $n$.
\end{lemma}

\begin{proof}
Consider the distinguished triangles
$$
K \otimes_\Lambda^\mathbf{L} \underline{I^n/I^{n + 1}} \to
K_{n + 1} \to
K_n \to
K \otimes_\Lambda^\mathbf{L} \underline{I^n/I^{n + 1}}[1]
$$
and the isomorphisms
$$
K \otimes_\Lambda^\mathbf{L} \underline{I^n/I^{n + 1}} =
K_1 \otimes_{\Lambda/I}^\mathbf{L} \underline{I^n/I^{n + 1}}
$$
By Lemma \ref{lemma-tensor-c} we see that this tensor product has
constructible cohomology sheaves (and vanishing when $K_1$ has
vanishing cohomology). Hence by induction on $n$ using
Lemma \ref{lemma-constructible-serre}
we see that each $K_n$ has constructible cohomology sheaves.
\end{proof}

\end{document}
